% =============================================================================
% Kapitel 1: Einleitung
% =============================================================================
\chapter{Einleitung}
\label{ch:einleitung}

\todo{Einleitenden Text verfassen, der das Thema der Arbeit vorstellt und den Leser in die Thematik einführt.}

% -----------------------------------------------------------------------------
\section{Motivation}
\label{sec:motivation}

\todo{Warum ist das Thema relevant? Welches Problem wird adressiert? Was ist der aktuelle Stand und wo besteht Handlungsbedarf?}

% -----------------------------------------------------------------------------
\section{Problemstellung}
\label{sec:problemstellung}

\todo{Konkrete Problemstellung formulieren. Was genau soll untersucht bzw.\ gelöst werden?}

% -----------------------------------------------------------------------------
\section{Zielsetzung}
\label{sec:zielsetzung}

\todo{Welche Ziele verfolgt diese Arbeit? Was soll am Ende erreicht werden?}

% -----------------------------------------------------------------------------
\section{Forschungsfragen}
\label{sec:forschungsfragen}

\todo{Die zentralen Forschungsfragen formulieren:}

\begin{enumerate}
    \item \todo{Forschungsfrage 1}
    \item \todo{Forschungsfrage 2}
    \item \todo{Forschungsfrage 3}
\end{enumerate}

% -----------------------------------------------------------------------------
\section{Aufbau der Arbeit}
\label{sec:aufbau}

\todo{Kurze Beschreibung der Struktur der Arbeit. Welches Kapitel behandelt welches Thema?}

Die vorliegende Arbeit ist wie folgt aufgebaut:

\textbf{\Cref{ch:grundlagen}} erläutert die theoretischen Grundlagen, die für das Verständnis dieser Arbeit notwendig sind.

\textbf{\Cref{ch:methodik}} beschreibt die gewählte Methodik und das Vorgehen.

\textbf{\Cref{ch:implementierung}} stellt die Implementierung der entwickelten Lösung vor.

\textbf{\Cref{ch:evaluation}} präsentiert die Evaluation und die erzielten Ergebnisse.

\textbf{\Cref{ch:diskussion}} diskutiert die Ergebnisse und deren Bedeutung.

\textbf{\Cref{ch:fazit}} fasst die Arbeit zusammen und gibt einen Ausblick auf zukünftige Arbeiten.
